\section{Number-theoretic algorithms: exact analytic gates}
\label{sec:number-algorithms}

Write \(\mathsf D\) for RH for all rational Dirichlet \(L\)-functions,
including \(\zeta_{\mathbb Q}\), and \(\mathsf H\) for the asserted RH for
Hecke \(L\)-functions over every number field. All character applications audited here are finite-order ones. The stipulated
Alp\"oge theorem discharges \(\mathsf D\). Every consequence of
\(\mathsf H\) below remains \emph{conditional on the all-number-field
claim of the Masullo draft}. A statement that a published proof's gate is
not discharged by \(\mathsf D\) is a scope audit, not a theorem of logical
nonimplication and not an exclusion of a different algorithm.

Finite fields are explicitly represented, for example as
\(\mathbb F_{p^m}=\mathbb F_p[T]/(g)\), with \(g\) supplied and irreducible.
A multiplicative primitive root generates \(\mathbb F_{p^m}^{\times}\);
a primitive element of the extension need only generate the field.
Integer factorizations supplied as input are stated explicitly.

\subsection{CX-N01: Miller testing with Bach's cutoff}
\label{cx-n01}

Handle bounded small integers and even inputs directly. For an odd
integer \(N\), write \(N-1=2^s u\), \(u\) odd, and test all
\[
 2\leq a\leq\min\{N-1,\lceil2(\log N)^2\rceil\}.
\]
Reject on a nontrivial gcd, or unless \(a^u\equiv1\pmod N\) or
\(a^{2^ju}\equiv-1\pmod N\) for some \(0\leq j<s\).
Under \(\mathsf D\), this is a deterministic polynomial-time primality test.

\paragraph{Exact source gate.}
Bach explicitly restricts \S6, pp.~369--373, to ordinary Dirichlet
characters; Theorem~2, p.~372, gives the subgroup bound with constant \(2\),
and p.~373 states the Miller/Solovay--Strassen witness consequence
\cite{CXNBach90}. Schoof's complete Theorem~2.2 proof, pp.~105--106,
explicitly states the smaller sufficient hypothesis of RH for
\emph{quadratic} Dirichlet characters \cite{CXNSchoof08}.

\paragraph{Full transfer proof.}
Those quadratic characters belong to \(\mathsf D\), so the cited criterion
gives a rejecting small base for every composite input. Every prime passes
all its nonzero bases. Increasing the strict source cutoff to the ceiling
cannot cause false acceptance, and bounded small inputs handle endpoint
conventions. There are \(O(b^2)\) bases for \(b=\lceil\log_2N\rceil\), each
requiring \(O(b)\) modular multiplications. The bit cost is
\(O(b^3M(b))\), where \(M(b)\) is multiplication cost; schoolbook arithmetic
gives \(O(b^5)\). Only the power of two in \(N-1\) is extracted.
This establishes this particular implementation; \(\mathrm{PRIMES}\)
was already in deterministic \(\mathrm P\) by AKS \cite[Theorem~4.1 p.~785 and Theorem~5.1 p.~789]{CXNAKS04}. \(\square\)

The entire Schoof proof was read for the hypothesis audit. Its squarefree
case uses quadratic characters modulo prime divisors and their products;
the repeated-prime case uses a smooth-number estimate and finite-group
counting. Miller's original locators are Theorem~2, p.~301, and the
Appendix, pp.~315--316 \cite{CXNMiller76}.

\subsection{CX-N02: a specified prime-field power nonresidue}
\label{cx-n02}

Given prime \(p\) and a specified prime \(r\mid p-1\), \(\mathsf D\) yields
an \(r\)-th power nonresidue in deterministic time polynomial in
\(\log p\), without factoring the rest of \(p-1\).

\paragraph{Proof.}
The subgroup \(G=(\mathbb F_p^\times)^r\) is proper.
If \(p\leq2(\log p)^2\), enumerate all nonzero residues.
Otherwise Bach's Theorem~2 supplies an integer outside \(G\) below that
cutoff, which is then below \(p\). In either case test candidates by
\(a^{(p-1)/r}\ne1\). Writing \(a=g^v\) in the cyclic group shows that this
is equivalent to \(r\nmid v\), exactly the desired condition.
The search has \(O((\log p)^2)\) candidates, with polynomial bit cost per
test. The analytic character with kernel \(G\) need not be computed.
\(\square\)

Bach's Theorem~2 quantifies over all positive integers below its bound.
Theorem~3, also p.~372, quantifies only over integers coprime to the modulus
and has constant \(3\). The preceding argument handles the possible zero
residue and does not conflate these statements. A fallback nonspecific
constant follows from Shoup's ordinary-character Proposition~9,
author-preprint p.~10 \cite{CXNShoup92}.

\subsection{CX-N03: certified primitive roots with factored order}
\label{cx-n03}

Given \(p\) and the complete prime factorization
\(p-1=\prod_{i=1}^t r_i^{e_i}\), a multiplicative generator of
\(\mathbb F_p^\times\) can be constructed in deterministic polynomial time
on \(\mathsf D\).

\paragraph{Proof.}
Use CX-N02 to find \(a_i\) with \(a_i^{(p-1)/r_i}\ne1\), and set
\[
              b_i=a_i^{(p-1)/r_i^{e_i}},\qquad b=\prod_i b_i .
\]
If \(a_i=g^{v_i}\), then \(r_i\nmid v_i\), so \(b_i\) has exact order
\(r_i^{e_i}\). In an abelian group, elements of pairwise coprime orders
have product of that product order: projection to every primary component
shows that any annihilating exponent is divisible by every component
order. Thus \(b\) has order \(p-1\). Since \(t\leq\log_2(p-1)\), the cost
is polynomial. The supplied factorization also certifies any candidate
through the tests \(b^{(p-1)/r_i}\ne1\). \(\square\)

No deterministic polynomial-time integer-factorization theorem is
provided here. A short guaranteed candidate list and efficient
identification of a generator are different outputs.

\subsection{CX-N04: Shoup's least primitive root bound}
\label{cx-n04}

For odd prime \(p\), Shoup's Theorem~3 gives the least positive primitive root modulo \(p\) as
\[
 O\!\left(t^4(\log t+1)^4(\log p)^2\right),\qquad t=\omega(p-1).
\]
The exact author-preprint locators are Theorem~3, p.~4; its proof in
\S6, p.~13; and Proposition~9, p.~10 \cite{CXNShoup92}.
All analytic characters used here are rational Dirichlet characters
modulo \(p\).

\paragraph{Proof of transfer and output.}
Apply \(\mathsf D\) to the finite family of character powers in the proof.
The remaining finite sieve is unconditional, so the bound and its
effective constants remain valid. The product of \(t\) different primes
is at least \(t!\); hence \(t(\log t+1)=O(\log p)\). The displayed bound is
therefore \(O((\log p)^6)\). Enumeration gives a polynomial-size list
containing a primitive root. Supplied factorization of \(p-1\) makes
recognition polynomial by CX-N03. The small bound alone does not provide
that factorization or a replacement recognition algorithm. \(\square\)

\subsection{CX-N05: Shoup's quadratic-extension candidate construction}
\label{cx-n05}

For odd prime \(p\), Shoup's different Theorem~2 constructs a candidate set for a primitive
root of \(\mathbb F_{p^2}\): author-preprint p.~3, proof \S5, pp.~12--13,
and Proposition~13 with the ray-class definitions in \S4, pp.~10--12
\cite{CXNShoup92}. Choose a prime \(\delta\equiv1\pmod4\) for which
\(-\delta\) is a quadratic nonresidue modulo \(p\), and set
\[
 K=\mathbb Q(\sqrt{-\delta}),\quad
 \mathcal O_K=\mathbb Z[\sqrt{-\delta}],\quad
 \mathfrak m=p\mathcal O_K .
\]
Proposition~10 uses rational characters to give
\(\delta=O((\log p)^2)\). Inertness identifies
\(\mathcal O_K/\mathfrak m\) with \(\mathbb F_{p^2}\).
The later character sums concern finite-order characters of
\(I_{\mathfrak m}/P_{\mathfrak m}^{+}\). Positivity is vacuous because
\(K\) is imaginary. These are Hecke characters over \(K\), together with
the trivial factor \(\zeta_K\). RH for the ray class field zeta
\(\zeta_{H_{K,\mathfrak m}}\) is sufficient.

\paragraph{Full transfer.}
Conditional \(\mathsf H\) covers every required character, so the
published character-sum and sieve estimates remain valid. The enumeration
is polynomial in \(\log p\), using the stated bounds on \(\delta\) and
the class number. Reduction gives the candidate set. An unconditional
deterministic field-isomorphism algorithm transfers it to a supplied
model \cite[Theorem~(1.2), p.~330]{CXNLenstra91}. With the factorization of \(p^2-1\), order tests identify a
generator; with only a specified prime \(r\mid p^2-1\), a power test
identifies an \(r\)-nonresidue in the list. \(\square\)

Here \(\mathsf D\) gives the initial prime selection and \(\zeta_K\),
because \(K/\mathbb Q\) is quadratic, but does not by this argument
supply the ray-class factors. The ray class field need not be abelian
over \(\mathbb Q\). This concerns the stated \emph{primitive-root candidate}
proof. It is not a necessity claim for each weaker task:
CX-N13 gives a smaller-hypothesis replacement for every specified
prime \(r\mid p-1\) in all explicitly given \(\mathbb F_{p^m}\),
including quadratic nonresidues.

\subsection{CX-N11: an explicit auxiliary Kummer field}
\label{cx-n11}

Let \(p\) be prime, \(n>1\), and \(p\nmid n\). Set
\[
              F=\mathbb Q(\zeta_{2n}),\qquad E=F(p^{1/n}).
\]
RH for \(\zeta_E\) and \(\zeta_{\mathbb Q}\) suffices for a prime
\[
 q=O(n^4\log^2(np)),\quad q\equiv1\pmod{2n},\quad
 \gcd\!\left((q-1)/\operatorname{ord}_q(p),n\right)=1.             \tag{N11.1}
\]
A search for it has polynomial cost in \(n,\log p\).
This supplies Buchmann--Shoup's Fact~2.1, p.~1314, and the closely
related Adleman--Lenstra Proposition~3, p.~351
\cite{CXNBuchmannShoup96,CXNAdlemanLenstra86}.

\paragraph{Exact analytic input.}
Bach--Sorenson's Theorem~5.1, p.~1729, gives an unramified rational prime
with any prescribed Frobenius conjugacy class in \(E/\mathbb Q\),
\(E\ne\mathbb Q\), satisfying
\[
              q\leq(4\log|\Delta_E|+2.5[E:\mathbb Q]+5)^2.      \tag{N11.2}
\]
Their convention in \S2, p.~1720 after (2.10), assumes RH for
\(\zeta_E\) and \(\zeta_{\mathbb Q}\) \cite{CXNBachSorenson96}.
There is no assumption of general Artin holomorphy.

\paragraph{Full proof of the field, conjugacy, and power conditions.}
The prime \(p\) is unramified in \(F\).
For odd \(n\), \(F=\mathbb Q(\zeta_n)\), which includes the case \(p=2\).
For even \(n\), \(p\) is odd and coprime to \(2n\).
Thus \(v_{\mathfrak P}(p)=1\) at every \(\mathfrak P\mid p\).
The polynomial \(X^n-p\) is Eisenstein over the discrete valuation ring
\(\mathcal O_{F,\mathfrak P}\), hence irreducible over \(F\).
Writing \(\alpha=p^{1/n}\), we obtain \([E:F]=n\).
The field \(E/\mathbb Q\) is Galois, since it contains all
\(\zeta_n^j\alpha\) and the \(2n\)-th roots of unity.

There is an \(F\)-automorphism \(\sigma(\alpha)=\zeta_n\alpha\) of order
\(n\). If \(\tau(\zeta_n)=\zeta_n^u\), then
\(\tau\sigma\tau^{-1}=\sigma^u\), as checked on \(F,\alpha\);
indeed \(\tau(\alpha)=\zeta_n^b\alpha\) for some \(b\).
The exponent \(u\) is a unit modulo \(n\). Thus each member of the
conjugacy class of \(\sigma\) fixes \(F\) and generates
\(\operatorname{Gal}(E/F)\).

Apply (N11.2) to this class. The resulting \(q\) splits in \(F\).
For even \(n\), whose cyclotomic conductor is \(2n\), this gives
\(q\equiv1\pmod{2n}\).
For odd \(n>1\), splitting gives \(q\equiv1\pmod n\), and \(q\ne2\)
since \(2\not\equiv1\pmod n\). As \(q\) is odd, it is again
\(1\pmod{2n}\). Also \(q\nmid pn\): the prime \(p\) ramifies in \(E/F\)
by the Eisenstein construction, and primes dividing \(n\) ramify in \(F\).

Choose a prime \(\mathfrak Q\mid q\) whose Frobenius is \(\sigma^u\).
In its residue field,
\(\alpha^q=\zeta_n^u\alpha\), so
\[
                    p^{(q-1)/n}=\alpha^{q-1}=\zeta_n^u .
\]
The reduction of \(\zeta_n\) has order \(n\), because \(q\nmid n\).
For every prime \(r\mid n\), therefore,
\[
                     p^{(q-1)/r}=(\zeta_n^u)^{n/r}\ne1.
\]
Writing \(p=g^a\) in \(\mathbb F_q^\times\), this says \(r\nmid a\).
Since \(\operatorname{ord}_q(p)=(q-1)/\gcd(a,q-1)\), it proves the gcd
condition in (N11.1).

Finally \(|\Delta_F|\leq(2n)^{[F:\mathbb Q]}\), and the relative
discriminant divides the polynomial discriminant ideal
\((n^np^{n-1})\). The discriminant tower formula gives
\[
 \begin{split}
 \log|\Delta_E|
 &\leq n[F:\mathbb Q]\log(2n)
       +[F:\mathbb Q](n\log n+(n-1)\log p)\\
 &=O(n^2\log(np)),\qquad [E:\mathbb Q]\leq2n^2.
 \end{split}
\]
Equation (N11.2) gives (N11.1), with effective constants.
Enumeration up to this bound is polynomial in \(n,\log p\);
primality and order testing, even using trial division of the
polynomially bounded \(q-1\), retain that cost. \(\square\)

The unique degree-\(n\) field \(K\subset\mathbb Q(\zeta_q)\) is real
because \(n\mid(q-1)/2\). The image of \(p\) in the cyclic quotient
\(\operatorname{Gal}(K/\mathbb Q)\) has order \(n\) by (N11.1);
hence \(p\) is inert and \(\mathcal O_K/p\mathcal O_K\cong\mathbb F_{p^n}\).
Adleman--Lenstra's Theorem~1, p.~350, Algorithm~A, pp.~350--351,
and Lemma~6, p.~352, use this mechanism to construct a field model.
That task is different from factoring an arbitrary input polynomial.

The extension \(E/F\) is cyclic. Equivalently, its finite family of
Hecke characters over \(F\) supplies \(\zeta_E\). Although \(\mathsf D\)
covers \(\zeta_F\), the additional Kummer characters are not thereby
rational Dirichlet characters. Conditional \(\mathsf H\) covers them.

\subsection{CX-N06: Buchmann--Shoup generating sets}
\label{cx-n06}

Buchmann--Shoup construct a model of \(\mathbb F_{p^n}\) and a set
generating its multiplicative group in deterministic time
\[
                              (n\log p)^{O(n)}.                \tag{N06.1}
\]
The main statement is on pp.~1311--1312; p.~1314 explicitly defines
ERH as RH for Dedekind zeta functions of all number fields
\cite{CXNBuchmannShoup96}.
For \(p\nmid n\), their auxiliary field \(K\) is the real cyclic field of
CX-N11. Fact~4.3, p.~1319, invokes Bach's Theorem~4 for the ray class
group modulo \(p\mathcal O_K\). A concrete sufficient package is
\[
                   \zeta_{\mathbb Q},\quad\zeta_E,\quad
                   \zeta_{H_{K,p\mathcal O_K}} .
\]
The ray modulus has no added infinite places, as in the paper.

\paragraph{Full transfer.}
Conditional \(\mathsf H\) supplies these zeta hypotheses. CX-N11
discharges prime selection; CX-N07 discharges ray-class generation.
The remaining lattice and finite-field operations preserve (N06.1).
For \(n=1\) and \(p>2\), Bach's Theorem~3 gives a group-generating set of size \(O((\log p)^2)\); \(p=2,n=1\) is trivial. If \(p\mid n\), then \(p\leq n\); enumerating irreducible-polynomial candidates and then all field elements costs \(p^{O(n)}\leq n^{O(n)}\), still within the stated bound. If \(S\) generates \(\mathbb F_{p^n}^\times\), then
for any specified prime \(r\mid p^n-1\), some member lies outside the
proper subgroup of \(r\)-th powers. Scanning with the exponent test
finds it. With complete factorization of \(p^n-1\), isolate primary
orders and multiply as in CX-N03. A generating set need not contain
one primitive root. \(\square\)

CX-N13 supplies a stronger nonresidue consequence for primes \(r\mid p-1\)
when the field model is already supplied.
For fixed \(n\), the time is polynomial in \(\log p\); for variable
\(n\), the displayed bound is not polynomial in the usual input
parameters. Discharging GRH does not change that degree dependence.

\subsection{CX-N07: class-group generators and their exact class field}
\label{cx-n07}

Bach's class-group bound is
\[
        \operatorname{Cl}(\mathcal O_K)=
        \langle[\mathfrak p]:N\mathfrak p\leq
                             12(\log|\Delta_K|)^2\rangle .     \tag{N07.1}
\]
The original locators are Theorem~4, p.~375, and the corollary on p.~376
\cite{CXNBach90}. Its analytic family comprises \(\zeta_K\) and the
nontrivial ideal-class-character \(L_K(s,\chi)\), which are unramified
finite-order Hecke \(L\)-functions over \(K\).

\paragraph{Proof from the precise character theorem.}
If those prime classes generate a proper subgroup \(G\), take a
nontrivial character of the quotient
\(\operatorname{Cl}(\mathcal O_K)/G\). Its conductor is \(1\).
The first bound in Bach's Theorem~4 gives a prime with
\[
 \chi(\mathfrak p)\ne1,\qquad
 N\mathfrak p\leq3\log^2(|\Delta_K|^2)=12(\log|\Delta_K|)^2,
\]
contradicting triviality on \(G\). The omitted degree-one case has
trivial class group. \(\square\)

\paragraph{Full class-field transfer.}
For the Hilbert class field \(H_K\), reciprocity gives
\[
 \operatorname{Gal}(H_K/K)\cong\operatorname{Cl}(\mathcal O_K),
 \qquad
 \zeta_{H_K}(s)=
  \prod_{\chi\in\operatorname{Cl}(\mathcal O_K)^\wedge}L_K(s,\chi).
                                                                    \tag{N07.2}
\]
The trivial factor is \(\zeta_K\); the other factors are entire.
None has a pole in \(0<\Re s<1\). Thus an off-line zero of any factor
would be an off-line zero of the product. RH for \(\zeta_{H_K}\)
suffices for all these factors. Conditional \(\mathsf H\) supplies it,
and hence (N07.1).

For a ray modulus \(\mathfrak m\), the same argument uses the ray
class field \(H_{K,\mathfrak m}\). Extra Euler factors from imprimitive
characters have zeros only on \(\Re s=0\), so do not affect the argument.
\(\square\)

Belabas--Diaz y Diaz--Friedman's refined computed bound assumes RH
only for the nontrivial class-character factors: author-preprint
Theorem~1.1 p.~3; Theorem~2.1 p.~5; proof pp.~4--6
\cite{CXNBDF08}. Their introduction distinguishes this from Bach's use
of both these factors and \(\zeta_K\). Neither gate should simply be
replaced by RH for \(\zeta_K\).

Even if \(K/\mathbb Q\) is abelian, so that \(\mathsf D\) covers
\(\zeta_K\), \(H_K/\mathbb Q\) need not be abelian. The additional
class-character gate remains in this proof. The bound gives generators,
not their relations. Algorithmic enumeration must also specify the
maximal order and its arithmetic, or account for computing them;
a defining polynomial is not silently the same input.

\subsection{CX-N08: verification of class groups and units}
\label{cx-n08}

Let \(S\) contain all archimedean places and a finite factor base.
Write \(K_S\) for the \emph{group of \(S\)-units} and \(H\subset K_S\)
for the subgroup generated by sampled \(S\)-units. For the valuation map
\[
 \phi:H\longrightarrow\mathbb Z^{S_{\rm fin}},\qquad
 x\longmapsto(\operatorname{ord}_{\mathfrak p}x)_{\mathfrak p},
\]
the tentative units and class group are \(\ker\phi\) and
\(\operatorname{coker}\phi\). In particular, \(K_S\) is not a relation
lattice. The exact source is Lenstra \S5, equation (5.7), author-preprint
p.~25 \cite{CXNLenstra92}.

\paragraph{Full verification proof.}
Assume the factor-base classes generate, \(H\) contains all roots of
unity, and the tentative groups have the required finite-index and rank
properties. The exact sequence
\[
 0\longrightarrow\ker\phi\longrightarrow\mathcal O_K^\times
 \longrightarrow K_S/H\longrightarrow\operatorname{coker}\phi
 \longrightarrow\operatorname{Cl}(\mathcal O_K)\longrightarrow0
\]
implies
\[
                              h'R'=hR[K_S:H],                 \tag{N08.1}
\]
where primes denote tentative class number and regulator.
The unit index is \(R'/R\), and the finite class-group index is
\(h'/h\); the exact sequence multiplies them to \([K_S:H]\).
If certified bounds give \(a/2<hR<a\) and the computed product also
satisfies \(a/2<h'R'<a\), then
\[
                         1\leq[K_S:H]=(h'R')/(hR)<2.
\]
The index is an integer, hence \(1\), proving completeness of both
units and class group. \(\square\)

Conditional \(\mathsf H\) supplies the factor-base bound. Where the
chosen analytic estimator assumes RH for \(\zeta_K\), it supplies
that premise also. The index test itself is algebraic and also works
with independently certified generators and analytic bounds.
It gives no rate for collecting sufficient relations. Any additional
smoothness heuristic or search-complexity assumption in a chosen
algorithm remains separate.

\subsection{CX-N12: a rigorously conditional randomized class-group algorithm}
\label{cx-n12}

Biasse--Erukulangara's modified Hafner--McCurley algorithm supplies a
concrete computation theorem \cite{CXNBiasseErukulangara}. We transfer
the input explicitly required by its algorithm boxes: a negative
\emph{fundamental} discriminant \(-d\), giving the full group structure
of \(\operatorname{Cl}(\mathcal O_{\mathbb Q(\sqrt{-d})})\).
The Main Theorem and Conjecture~1 are on author-preprint p.~4,
Algorithm~4 on p.~15, and Corollary~1 on p.~16.
Conjecture~1 explicitly concerns Hecke \(L\)-functions over number fields.

Conditional on \(\mathsf H\), the algorithm has Las Vegas expected
bit complexity
\[
 \exp\!\left(\left(\frac3{\sqrt8}+o(1)\right)
                    \sqrt{\log d\,\log\log d}\right).           \tag{N12.1}
\]
It remains randomized. The displayed bound is superpolynomial and
subexponential in the bit length \(\log d\).

\paragraph{The visible number-field family.}
The expansion step, Proposition~1, author-preprint p.~8, uses
Jao--Miller--Venkatesan's Theorem~1.1 and Corollary~1.3
\cite[arXiv version, pp.~3--4; proof pp.~9--10]{CXNJMV09}.
For modulus \(1\) in imaginary quadratic \(K\), their characters are
exactly those of \(\operatorname{Cl}(\mathcal O_K)\), and their proof
explicitly estimates the corresponding Hecke \(L\)-functions.
Thus RH for \(\zeta_{H_K}\) supplies this step and class generation
by CX-N07. The rational quadratic character of \(K\) controls the
separate quadratic splitting and class-number estimates.
The fact that \(K/\mathbb Q\) is quadratic does not itself eliminate
the class-character gate.

For completeness, the split-prime factor base in the algorithm can be
obtained from the same family. Regard each class character as an
imprimitive character modulo \(d\mathcal O_K\), extended by zero at
primes dividing \(d\). The second bound of Bach's Theorem~4 gives
\(\chi(\mathfrak p)\notin\{0,1\}\) with
\[
 N\mathfrak p\leq12\log^2(d^2N(d\mathcal O_K))
              =192(\log d)^2.
\]
Such a prime is unramified; it cannot be an inert prime ideal, since
an inert prime ideal is the principal ideal \((p)\) and has class
character value \(1\). It therefore lies over a split rational prime.
Separating a proper generated subgroup by a character proves that
these split-prime classes generate. Choosing only one of the two
primes above each split rational prime preserves generation, because
the two classes are inverses. Only finite Euler factors have
been added, so no additional number-field character family is needed.

\paragraph{Full transfer, certification, and expected time.}
Conditional \(\mathsf H\) supplies the precise hypothesis explicitly
declared in Conjecture~1. Hence the paper's conditional factor-base,
relation-probability, and class-number estimates apply with the
published parameters.

Let \(\Lambda\subset\mathbb Z^n\) be the full relation lattice for the
factor base, and \(\Lambda'\subset\Lambda\) the collected full-rank
lattice. Every collected vector is a true relation. Since the factor
base generates,
\[
       h'=\det\Lambda'=h[\Lambda:\Lambda'],\qquad
       h=|\operatorname{Cl}(\mathcal O_K)|.
\]
Algorithm~4, steps 15--20, computes a certified \(h^*\) with
\(h^*\leq h<2h^*\), obtains \(h'\) from the Smith normal form, and
declares \textsc{Failure} if \(h'\geq2h^*\).
An accepted output satisfies \(1\leq h'/h<2h^*/h\leq2\);
the integer index is therefore \(1\). Every accepted group is correct
under the analytic estimates, and unsuccessful branches are declared
failures. The Smith-form procedure is itself Las Vegas.

Put \(L(d)=\exp(\sqrt{\log d\,\log\log d})\).
Corollary~1 supplies a time-bounded run of cost
\(L(d)^{3/\sqrt8+o(1)}\) with success probability at least
\(1-1/d^{1+o(1)}\). Independent restarts on failure have a geometric
number of trials, with expectation at most \(2\) for sufficiently
large \(d\); the bounded smaller inputs can be handled directly.
This gives (N12.1) with its randomness retained. The corollary balances
\[
             \min_{z>0}\max\{3z,z+1/(4z)\}=3/\sqrt8,
             \qquad z=1/\sqrt8.
\]
The expectation deduction is justified by detectable failures and
correctness of accepted outputs, not by high success probability alone.
\(\square\)

\subsection{CX-N09: general finite-field factorization}
\label{cx-n09}

Evdokimov's published theorem gives deterministic factorization of a
degree-\(d\) polynomial over an explicit field \(\mathbb F_q\) in
\[
                            (d^{\log d}\log q)^{O(1)}            \tag{N09.1}
\]
under the relevant GRH. The primary statement and radical-oracle
description are in the abstract on p.~209 \cite{CXNEvdokimov94}.
The degree dependence is quasipolynomial. For fixed \(d\), the bound
is polynomial in \(\log q\); it is not a general polynomial bound
in \(d,\log q\). Unconditional randomized polynomial-time factorization
already predates this result \cite[introduction, author pp.~1--2]{CXNGao01}.

\paragraph{A sufficient broad radical gate.}
Evdokimov's earlier Theorem~I.3, pp.~109--110, gives deterministic
\(k\)-th root extraction in an explicitly given \(\mathbb F_{p^m}\)
in time polynomial in \(k,m,\log p\) \cite{CXNEvdokimov89}.
Its proof uses field construction and isomorphisms from
Theorems~I.1 and I.2, pp.~107--109. The analytic convention on p.~105
explicitly allows \emph{Dirichlet characters of abelian number fields}.
These are finite-order Hecke characters over those fields.
For example, case III of Theorem~I.1, p.~108, uses the quadratic
residue character modulo \(p\) over \(K=\mathbb Q(i)\), identifying
its ideal character using the fact that \(\mathbb Z[i]\) is a PID.
The hypothesis is not automatically the rational Dirichlet assertion.

Every character allowed by that oracle theorem belongs to
\(\mathsf H\). Conditional on \(\mathsf H\), the radical oracle has
the required polynomial bound. Substituting it into the published
radical-oracle reduction gives (N09.1), with its degree dependence
unchanged. This uses no assumption of arbitrary Artin holomorphy.
It is a transfer of the cited algorithmic theorem, not an independent
reproof of the complete Evdokimov algebraic algorithm.
CX-N15 supplies an independent composition of the same bound from
accessible primary algorithm statements, with explicit Kummer-zeta
hypotheses and without that original internal reduction.

\paragraph{Why a smaller-hypothesis transfer needs its own audit.}
Gao proves that radicals in an elementary algebra
\(\mathbb F_q^N\) require only \(r\)-nonresidues in the base
\(\mathbb F_q\), for primes \(r\mid\gcd(k,q-1)\):
\S2.4, author-preprint pp.~7--9, Lemma~2.6 \cite{CXNGao01}.
His introduction, pp.~1--2, records the unconditional reduction to
completely split prime-field polynomials.
These facts make a rational-character route a legitimate question;
they do not alone audit subsequent automorphism splitting.
Gao's Lemma~3.6, author-preprint p.~16, invokes the latter from
R\'onyai's work.

The broad \(\mathsf H\)-transfer is therefore established here, while
a \(\mathsf D\)-only transfer of the entire general theorem is not
asserted. The separately proved CX-N14 handles several specified
automorphism-order cases, including prime-order powers whose order divides
\(p-1\), using at most rational characters. The unrestricted step
remains a stated gate in this audit.

\subsection{CX-N10: special factorization families preserve their parameters}
\label{cx-n10}

Bach--von zur Gathen--Lenstra's Theorem~1, p.~6, is a precise
special-field theorem \cite{CXNBGL01}. Let \(S(p^k)\) be the prime
divisors of the integer \(\Phi_k(p)\), \(s(p^k)\) its largest prime
divisor, and \(R(p^k)\) the primes dividing \(\ell-1\) for some
\(\ell\in S(p^k)\). For an explicit \(\mathbb F_{p^n}\) and \(f\),
the theorem assumes a degree-\(r\) irreducible polynomial over
\(\mathbb F_p\) for each \(r\in R(p^k)\) not dividing \(n\), and has time
\[
                     (s(p^k)+\deg f+n\log p)^{O(1)}.            \tag{N10.1}
\]

\paragraph{Transfer.}
Adleman--Lenstra constructs the extra irreducible-polynomial inputs
when its auxiliary-prime hypothesis is supplied.
For each required prime \(r\ne p\), CX-N11 supplies the concrete
field \(\mathbb Q(\zeta_{2r},p^{1/r})\); characteristic-\(p\)
extensions use the separate algebraic construction.
Conditional \(\mathsf H\) covers these hypotheses and removes the
extra input gates while retaining (N10.1).
The same paper's Theorem~2, p.~7, constructs a primitive
\(\ell\)-th root of unity from specified irreducible-polynomial inputs;
Theorem~3, p.~7, constructs the relevant nonresidues from roots of
unity. These intermediate tasks have explicit inputs and parameters.
They do not give a polynomial bound for arbitrary polynomial
factorization. \(\square\)

For Fermat or Mersenne characteristics, \(k=1\) or \(k=2\) makes
\(\Phi_k(p)\) a power of \(2\), leaving \(R(p^k)\) empty.
The corollary on p.~6 is then unconditional, rather than newly
unconditional by either GRH claim.

R\'onyai's other results also have structural parameters.
The primary abstract of \cite{CXNRonyai92} gives time polynomial in
the input degree, coefficient size, \(\log p\), and the rational
splitting-field degree; the latter may be exponential in the usual
input length. Gao's introduction, p.~2, records R\'onyai's 1988
algorithm for finding a nontrivial factor of a degree-\(d\) completely
split polynomial, with cost polynomial in \(d^r,\log p\), for a
prime divisor \(r\mid d\). It is polynomial for bounded \(r\), not
uniformly when \(r\) grows. These are sourced parameter distinctions.
The full original proofs and precise theorem numbers of these two
R\'onyai statements were not independently retrieved; no stronger
rational-Dirichlet transfer is inferred from the acronym GRH.
